\subsection{Retained scalar and partial-swap contract}
The bit payload is $\mathbb F_2$; its address geometry is rational and
is later reduced at a suitable odd prime. Let $v_h=\binom h3$ and use
$H_h=I-J/9$ on triple indicators. Their inner product is
$|S\cap T|-1$, so intersection-one lines are orthogonal and each line
has norm two. The retained paired exclusion producer computes the
side sums and the $h$ point totals. Its scalar copy, reversible mixer,
and scatter maps satisfy $JLV=I$. The chronological dirty-scratch word
\[
 L,\ J,\ L^{-1},\ V,\ L,\ J,\ L^{-1},\ V
\]
adds $JLz+JL(z+Vx)=x$ and restores every arbitrary auxiliary value $z$.
Three alternating data shears implement exchange. The positive source
envelopes and selected compatible enlargements make all required
projector differences nondegenerate. Carrier reuse joins only nested,
properly ordered continuations; it changes physical roles, not the
scalar map. These are the construction and histogram identities of
\path{notes/endpoint-gauge-construction.tex}.

For a rational projector $P$ the partial swap is
\[
 D_P=\begin{pmatrix}I-P&P\\P&I-P\end{pmatrix},\qquad D_P^2=I.
\]
If $U\subset V$ are nondegenerate then
$D_{P_V}D_{P_U}=D_{P_V-P_U}$. The retained exact lower-triangular
conjugation sends $D_P$ to its cross-bank pivot swaps: if
$P=L\Pi R$, put $C=RL$, $Q_X=\Pi\Pi^t$, $Q_Y=\Pi^t\Pi$, and
\[
 K_0=Q_YC-CQ_X,\qquad
 G_0=\begin{pmatrix}L&0\\0&R^{-1}\end{pmatrix}
     \begin{pmatrix}I&0\\K_0&I\end{pmatrix}.
\]
Then
\[
 G_0^{-1}D_PG_0=
 \begin{pmatrix}I-Q_X&\Pi\\\Pi^t&I-Q_Y\end{pmatrix}.
\]
All wrappers are ordered-affine operations. A contiguous diagonal run
of $t$ pivots therefore uses one width-$t$ interchange, with no gathering
of separated positions. The address prime excludes the finitely many
bad denominators and nonzero minors after the rational basis is fixed.

For stage two use $P_0=(I-P_a)\otimes I_b\otimes P_k$,
$Q=P_a\otimes I_b\otimes P_k$ and $P_1=P_0+Q$.
The source $D_{P_0}$ and sink $D_{I-P_0}$ have endpoint product $D_I$
on every restored auxiliary role. The retained inverse-copy schedule
raises source data to $P_0$ before copying, and all subsequent producer
frames contain $P_0$. The entrance vanishes; the exit is
$D_{I-P_0}D_{P_1}=D_{I-Q}$ of rank $m-b$. Its rank equals the sum of
its two former boundary charges. This applies to the full stage-two
bank, including fanout and reused carriers.

\subsection{One controlled basis with a fixed middle factor}
Fix $(a,b,k)=(32,30,36)$, $H=ab=960$, and $m_b=Hk=34560$.
Use $i=b\alpha+\beta$ as inner physical index and
\[
 K=T_2(I_a\otimes L_b),\quad
 T_2(v\otimes e_\beta)=M_\beta L_av\otimes e_\beta,
 \quad L_b=I_b+J_b.
\]
Keep the inherited inner boundary prescription: the first $2a$ row
labels are $[I_a,\operatorname{cyc}_{+1}I_a]$ and the last $2a$
inverse-column labels are $[\operatorname{cyc}_{-1}I_a,I_a]$.
A label $c$ at $(\alpha,\beta)$ prescribes
$e_\alpha^tM_\beta=e_c^t$ or $M_\beta^{-1}e_\alpha=e_c$,
respectively. The residue sets $\{0,1,3\}$ and $\{25,27,28\}$ are
disjoint. Thus each $M_\beta$ has a permutation completion, and the
remaining invertible locus is a nonempty irreducible rational family.
The first and third axis bases $L_a,L_k$ remain generic.

Set $S=T(I_k\otimes K)$ with
$T(u\otimes e_i)=G_i u\otimes e_i$. Match only the first $k$ rows
of $G_i$ to the coordinate rows of $L_k$, and the last $k$ inverse
columns to the matching columns of $L_k^{-1}$. The next $a$ rows
and preceding $a$ inverse columns are free. The two length-$(a+k)$
boundary intervals are disjoint.

Fixing $L_b$ requires a new nonvanishing justification. For each triple
indicator $t$,
\[
 L_bt=t+3\mathbf1,\qquad
 \phi_tL_b^{-1}=\tfrac12\left(t-\frac{10}{3(b+1)}\mathbf1\right)^t.
\]
Every primal and dual coordinate is nonzero. Distinct triple lines
remain independent, so their coordinate-ratio vector is nonconstant.
These facts preserve the diagonal stage-two corner, the inner tree
witness, and the matched outer A3 corner.

For completeness, the A1 null projector is
\[
 Q_1=I_k\otimes P_q+P_t\otimes P_U,\qquad
 \operatorname{rank}P_q=1,\quad\operatorname{rank}P_U=a,\quad P_qP_U=0.
\]
After normalizing line coordinates its first $k$ nullspace rows have
form $[e_p^t,t_pz_p]$, where $t_p\ne0$. The first $a$ points $z_p$
are nonzero multiples of a coordinate basis and span an affine
hyperplane $\mathcal H$. The repeated labels among the first $2a$
points compare adjacent fast-coordinate ratios around the entire
$b$-cycle. If all these points lay in $\mathcal H$, every such ratio
would agree, contradicting independence of the two motif lines.
Since $2a\le a+k$, either the first $k$ points already affinely span
$\mathbb Q^a$, or an extra point lies outside $\mathcal H$.
For a free extra row $g_i$, elimination of the first $k$ coordinates
leaves
\[
 \sum_{p<k}g_{i,p}t_p(z_i-z_p).
\]
In the first case this map is onto $\mathbb Q^a$ for every extra row.
In the second it contains the $(a-1)$-dimensional direction space of
$\mathcal H$ for every row, and the transverse row supplies the last
direction. Choose the $a$ rows to form a basis. The reverse cyclic
inverse-column prescription gives the dual witness. Thus both full
$(a+k)$-dimensional nullspace restrictions are nonsingular in this
restricted family.

The A5 nullity is $a+b-1$. After the first $b$ eliminations, the two
restrictions are incidence matrices on vertices $0,\ldots,b+1$ with
edges, respectively,
\[
 (0,b),(1,b+1),\quad(j+2,j+1)\ (0\le j\le b-3),\quad(0,b-1),
\]
\[
 (b,0),(b+1,1),(b+1,2),\quad(j-2,j-1)\ (4\le j\le b+1).
\]
Both are trees and have rank $a-1$ by leaf elimination. Their diagonal
rescalings use only the nonzero evaluations displayed above. The A3
null projector $P_t\otimes I_H+(I-P_t)\otimes P_q$ has full nullity
$H+k-1$: eliminate its first $H$ diagonal restrictions, then use the
free second rows of $G_i$ and free next-to-last inverse columns to
span the $(k-1)$-dimensional complement of $t$. Its matched $k$-corner
is the ordinary corank-one projector, yielding one singleton and a
$(k-2)$-block; the unperturbed interval supplies $H-2k+2$.
These are precisely the corner arguments of
\path{notes/endpoint-gauge-basis.tex}, with the fixed-middle
nonvanishing and A1 cycle argument supplied here.

All these determinants, and the generic first/third-axis internal
corner determinants, are nonzero on the same remaining irreducible
family. Their finite product is nonzero. Rational enumeration provides
one simultaneous basis for the whole finite circuit. Middle-axis
internal profiles instead use the exact certificate in the next section.

\subsection{The refined A1 and A5 contiguous blocks}
In the first-$a$/last-$a$ A1 corner, the $P_t\otimes P_U$ term is an
invertible diagonal matrix. The $I_k\otimes P_q$ term can be nonzero
only at row $i=(k-a)+j=4+j$, a strictly lower diagonal. The corner is
therefore invertible lower triangular and can be normalized by one
allowed lower row operation. It gives a single $a$-block. The remaining
$k$ pivots in the full nullity corner are retained as singletons.
The large-projector Schur identity leaves $m_b-2(a+k)$, so A1 gives
36 singletons and blocks 32 and 34424.

For A5, the first-$b$/last-$b$ corner has form $B-uv^t$, up to
nonzero diagonal scales, where $B$ is lower triangular with nonzero
diagonal. Its two contributions are the diagonal $P_a\otimes I_b$
and the offset-two lower diagonal $I_a\otimes P_b$; the remaining
term is rank one. Since $u_0v_{b-1}\ne0$, pivot at the extreme
upper-right entry. On indices $1,\ldots,b-2$ its Schur complement is
exactly the corresponding block of $B$: the first row and last column
of $B$ have zero entries in this central interval, so elimination
cancels the rank-one term there. This gives one contiguous $(b-2)$-block.
The remaining $a$ pivots of the nonsingular full nullity corner remain
singletons. Including the extreme pivot and the large-projector middle
block yields 33 singletons and blocks 28 and 838.

\subsection{Counts and strict characteristic}
The selected per-axis counts are
\[
\begin{array}{c|rrrrrr}
h&v_h&c_h&q_h&|\mathcal M_h|&R_h&\text{loss}_h\\\hline
32&4960&105056&14912&3296&116672&992\\
30&4060&85440&12210&2910&94740&870\\
36&7140&153144&21456&4140&170460&1260
\end{array}
\]
The middle axis uses its original positive envelopes and physical
transition multiset; the former enlarged-positive rank histogram does
not determine its selected pivot runs. Each selected internal rank
mass is $hR_h+2h(h-1)$. Put $N=v_av_bv_k$, $B_i=NR_{h_i}/v_{h_i}$.
The retained intersection-one bipartite matching shares the first and
third auxiliary banks. Thus
\[
\begin{aligned}
N&=143782464000,& B_1&=3382134604800,\\
B_2&=3355160256000,& B_3&=3432655296000,\\
W_b=2N+B_2+B_3&=7075380480000,&L_b^{\rm loss}&=84940396800,\\
s_b=W_bm_b-2N+2L_b^{\rm loss}&=244525031704665600.
\end{aligned}
\]
Here $L_b^{\rm loss}$ is a scalar loss, distinct from the fixed basis.
The complete external list is
\[
\begin{array}{c|r|l}
\text{class}&\text{copies}&\text{widths}\\\hline
\text{A1}&B_1&36\text{ singletons};32;34424\\
\text{third-only}&B_3-B_1&36;34488\\
\text{gauged stage two}&B_2&30;34500\\
\text{A3}&2N&71\text{ singletons};34;890;32570\\
\text{A5}&2N&33\text{ singletons};28;838\\
\text{data fronts}&2N\text{ each}&(1,30);(1,28);(1,34)
\end{array}
\]
For $h=32,36$, each rank-$r$ internal histogram entry contributes
$N/v_h$ times its multiplicity: use $r$ singletons when $2r\le h$,
and $h-r$ singletons plus one width-$(2r-h)$ child otherwise.
For $h=30$ use the next section's certified block multiplicities,
replicated $N/v_{30}$ times. These disjoint classes define $n_t$,
with $\sum_ttn_t=s_b$ and $0<t\le34500<m_b$.

Set $a_b=1252373/10^{11}$ and $\tau=1-a_b$. To enclose logarithms,
for $1\le y\le2$ and $z=(y-1)/(y+1)$ use
\[
 \log y\le2\sum_{j=0}^{23}\frac{z^{2j+1}}{2j+1}
             +\frac{2z^{49}}{49(1-z^2)}.
\]
Power-of-two extraction and upward rounding to $10^{-12}$ give
$\ell_t\ge\log(m_b/t)$. The bound
$e^u\le1+u+u^2/[2(1-u/3)]$ for $0\le u<3$ yields
\[
 \Psi_b(\tau)\le\sum_t\frac{n_tt}{W_bm_b}
 \left(1+a_b\ell_t+\frac{(a_b\ell_t)^2}{2(1-a_b\ell_t/3)}\right)
 <1-2.8159\times10^{-14}.
\]
The final inequality is a strict rational comparison, not a numerical
root estimate. Strong induction with the inherited floor-and-remainder
construction proves the arbitrary-width interchange interface
$O(Ve^\tau)$, with complete spectators and restored scratch.
